% -*- coding: utf-8 -*-
% !TEX program = xelatex
\documentclass[plain,noanswer,medmath]{../../template/jnuexam}
\usepackage{../../template/kysx}

\begin{document}


\examtitle{2008}{4}


\loadproblems{1}{2008P1}%载入试卷一
\loadproblems{2}{2008P2}%载入试卷二
\loadproblems{3}{2008P3}%载入试卷三


\makepart{选择题}{1～8小题，每小题4分，共32分}


\begin{problem}
设$0 < a < b$，则$\lim_{n\to\infty}\left(a^{-n} + b^{-n}\right)^{\frac{1}{n}}=$\pickout{}
\begin{abcd}
\item $a$．
\item $a^{-1}$．
\item $b$．
\item $b^{-1}$．
\end{abcd}
\end{problem}


\useproblem{3}{1}{1}%【同试卷三第一(1)题】


\begin{problem}
设$f(x)$是连续的奇函数，$g(x)$是连续的偶函数，区域\goodbreak
$D = \{(x,y)|0 \le x \le 1, - \sqrt{x} \le y \le \sqrt{x}\}$，则以下结论正确的是\pickout{}
\begin{abcd}
\item $\iint_D f(y)g(x)\dx\dy = 0$．
\item $\iint_D f(x)g(y)\dx\dy = 0$．
\item $\iint_D [f(x) + g(y)]\dx\dy = 0$．
\item $\iint_D [f(y) + g(x)]\dx\dy = 0$．
\end{abcd}
\end{problem}


\useproblem{2}{1}{2}%【同试卷二第一(2)题】


\useproblem{1}{1}{5}%【同试卷一第一(5)题】


\useproblem{2}{1}{8}%【同试卷二第一(8)题】


\useproblem{1}{1}{7}%【同试卷一第一(7)题】


\useproblem{1}{1}{8}%【同试卷一第一(8)题】


\makepart{填空题}{9～14小题，每小题4分，共24分}


\useproblem{3}{2}{9}%【同试卷三第二(9)题】


\begin{problem}
已知函数$f(x)$连续，且$\lim_{x \to 0} \frac{f(x)}{x} = 2$，
则曲线$y = f(x)$上对应$x = 0$处的切线方程是 \fillin{}．
\end{problem}


\begin{problem}
$\int_1^2 \dx \int_0^1 x^y\ln x\dy =$ \fillin{}．
\end{problem}


\useproblem{2}{2}{10}%【同试卷二第二(10)题】


\begin{problem}
设$3$阶矩阵$A$的特征值互不相同，且行列式$|A| = 0$，则$A$的秩为 \fillin{}．
\end{problem}


\useproblem{1}{2}{14}%【同试卷一第二(14)题】


\makepart{解答题}{15～23小题，共94分}


\useproblem[points=9]{3}{3}{15}%【同试卷三第三(15)题】


\begin{problem}[points=10]
设函数$f(x) = \int_0^1 |t(t - x)|\dt$（$0 < x < 1$），
求$f(x)$的极值、单调区间及曲线$y = f(x)$的凹凸区间．
\end{problem}


\useproblem[points=11]{2}{3}{21}%【同试卷二第三(21)题】


\useproblem[points=10]{3}{3}{16}%【同试卷三第三(16)题】


\useproblem[points=10]{3}{3}{18}%【同试卷三第三(18)题】


\useproblem[points=12]{1}{3}{21}%【同试卷一第三(21)题】


\useproblem[points=10]{2}{3}{23}%【同试卷二第三(23)题】


\useproblem[points=11]{1}{3}{22}%【同试卷一第三(22)题】


\begin{problem}[points=11]
设某企业生产线上产品合格率为$0.96$，
不合格产品中只有$\frac{3}{4}$产品可进行再加工，且再加工合格率为$0.8$，其余均为废品．
每件合格品获利$80$元，每件废品亏损$20$元，
为保证该企业每天平均利润不低于$2$万元，问企业每天至少应生产多少件产品？
\end{problem}


\end{document}
